% ============================================================
%  paper-export :: 封面
%
%  pandoc 默认的 \maketitle 把标题/作者/日期挤在页面上部，日期紧跟作者，
%  既不像交付文档也不像论文。这里改成独立封面：
%    顶部留白 ≥15% 页高 → 单位名 → 主标题 → 副标题 → 分隔线
%    → 信息栏（文档编号/版本/密级/编制人）→ 弹性留白 → 日期 → 底部留白
%
%  字段值由 export.sh 生成的临时 cover-vars.tex 用 \def 注入（已做 LaTeX 转义）。
%  未提供的字段整行省略，不留空行。
% ============================================================

\makeatletter

% ---------- 字段默认值（export.sh 会覆盖）----------
\providecommand{\peorg}{}          % 编制单位，显示在标题上方
\providecommand{\pesubtitle}{}     % 副标题
\providecommand{\pedocno}{}        % 文档编号
\providecommand{\peversion}{}      % 版本号
\providecommand{\pesecurity}{}     % 密级
\providecommand{\peauthorline}{}   % 编制人
\providecommand{\pehastoc}{1}      % 是否生成目录（供页码分节判断）

% 字段是否为空。参数传的是宏本身（\pedocno 这样的单 token），
% 直接 \ifx 与 \@empty 比较即可，不必 \edef —— 值里可能含 \& \% 之类的
% 转义序列，\edef 展开它们会炸。
\newcommand{\pe@ifempty}[3]{\ifx#1\@empty#2\else#3\fi}

% 信息栏的一行。刻意不用 tabular：单元格分隔符 & 与换行 \\ 落在未闭合的
% \ifx 分支里会触发「Incomplete \ifx」—— 对齐扫描会打断条件扫描。
% 改用两个零宽盒子共用同一个锚点：标签向左伸、值向右伸。
% 信息栏的锚点比页面中线左移约 6mm，让整组字段在封面上更稳。
\newcommand{\pe@lbl}[1]{\makebox[4.8em][s]{#1}}
\newcommand{\pe@inforow}[2]{%
  \pe@ifempty{#2}{}{%
    \hspace*{-1.2em}%
    \makebox[0pt][r]{\pe@lbl{#1}}%
    \makebox[0pt][l]{\hspace{0.15em}：\hspace{0.35em}#2}%
    \par\vspace{0.5\baselineskip}}}

% ---------- 封面 ----------
\renewcommand{\maketitle}{%
  \begin{titlepage}
    % 正文为装订留出较宽左边距；封面仍应按纸张物理中心排版。
    \addtolength{\oddsidemargin}{\dimexpr(\Gm@rmargin-\Gm@lmargin)/2\relax}%
    \centering
    \color{pblack}
    % 顶部留白：主标题绝不贴顶
    \vspace*{0.15\textheight}

    % 编制单位（可选）
    \pe@ifempty{\peorg}{}{%
      {\heiti\zihao{3}\peorg\par}
      \vspace{0.055\textheight}}

    % 主标题：一号黑体，长标题自动折行且行距放宽
    {\heiti\zihao{1}\linespread{1.45}\selectfont\@title\par}

    % 副标题（可选）
    \pe@ifempty{\pesubtitle}{}{%
      \vspace{1.1\baselineskip}
      {\heiti\zihao{3}\linespread{1.35}\selectfont\pesubtitle\par}}

    % 分隔线
    \vspace{0.05\textheight}
    {\color{prule}\rule{0.58\textwidth}{0.7pt}}

    % 信息栏：文档编号 / 版本 / 密级 / 编制人
    \vspace{0.055\textheight}
    {\kaishu\zihao{-3}%
      \pe@inforow{文\hfill 档\hfill 编\hfill 号}{\pedocno}%
      \pe@inforow{版\hfill 本\hfill 号}{\peversion}%
      \pe@inforow{密\hfill 级}{\pesecurity}%
      \pe@inforow{编\hfill 制\hfill 人}{\peauthorline}%
      \par}

    % 弹性留白把日期压到页面下部
    \vfill
    {\songti\zihao{3}\@date\par}
    \vspace{0.10\textheight}
  \end{titlepage}%
  % 封面不参与页码计数：前置部分从 I 重新开始
  \setcounter{page}{1}%
}

\makeatother
